% SOURCE_AUTHORITY: sga5_fr_workpass.tex, lines 9045--9105 (LF-normalized unit).
% SOURCE_SHA256: 791F4EFFC5E02832D5D77ED03518C8156D6F07E4C8238B03545DB93D883FBB28
\section*{Exposição VII}

\begin{center}
\textbf{COHOMOLOGIA DE ALGUNOS ESQUEMAS CLÁSSICOS}\\[0.5em]
\textbf{E}\\[0.5em]
\textbf{TEORIA COHOMOLÓGICA das CLASSES DE CHERN}\\[1em]
por\\[0.5em]
J. P. JOUANOLOU
\end{center}

Em todo o texto supõe-se dado um inteiro $\nu>0$ e define-se
\[
A=\mathbb Z/\nu\mathbb Z;
\]
As características residuais de todos os esquemas considerados são primas com o inteiro $\nu$. Se $X$ é tal esquema, denota-se por $\mu_X$ o feixe das raízes $\nu$-ésimas da unidade e por $A^{\bullet}(X)$ o anel de cohomologia
\[
\bigoplus_i H^{2i}(X,\mu_X^{\otimes i})
\]
cujo produto é o produto cup.

\subsection*{1. Fibrados vetoriais}

Sejam $S$ um esquema e $E$ uma $\mathcal O_S$-módulo localmente livre, entendido de tipo finito. Define-se $X=\V(E)$, $U=\V(E)^*$, complementar à seção nula, e são denotados por
\[
p:X\to S,
\qquad
q:U\to S
\]
as projeções canônicas.

\begin{statement}{Proposição 1.1}
Seja $L$ um $A$-módulo.
\begin{enumerate}[label=(\roman*)]
\item
\begin{enumerate}[label=\alph*)]
\item Para todo inteiro $j>0$, temos
\[
\R^j p_*(p^*L)=0.
\]
\item O morfismo de adjunção
\[
L\longrightarrow p_*p^*(L)
\]
é um isomorfismo.
\end{enumerate}
\item Supõe-se que $E$ tem posto constante $r$.
\begin{enumerate}[label=\alph*)]
\item
\[
\R_!^i p(p^*L)=0\quad\text{se }i\ne 2r.
\]
\item O morfismo de traço (SGA 4 XVIII)
\[
\R_!^{2r}p(p^*L)\longrightarrow L\otimes\mu_S^{\otimes -r}
\]
é um isomorfismo.
\end{enumerate}
\end{enumerate}
\end{statement}

\emph{Demonstração.} A parte (i) segue-se imediatamente, por localização para a topologia de Zariski, de (SGA 4 XV 2.2). A afirmação (ii) está essencialmente contida em (SGA 4 XVIII).
